\section{Nonequilibrium Modeling}

\subsection{General Model}

In this section we derive a general nonequilibrium kinetic model used for fitting the MWS data and viability modeling. We consider in the rate of change of the electron density, \(n_{e}(t)\) in time. 

\begin{equation}
  \label{eq:start_deqn}
\frac{dn_{e}(t)}{dt} = G - R
\end{equation}

Where \(G\) is the generation rate of electrons and \(R\) is the recombination rate. We consider the generation rate to be comprised of a thermal generation rate and a nonequilibrium generation rate.

\begin{equation}
  G = G_{th} + G_{NE}
\end{equation}

Where \(G_{th}\) is the thermal generation rate constant and \(G_{NE}\) is the Nonequilibrium generation rate. 

The recombination rate can be written

\begin{equation}
  R = \sum_{i}^{}k_{r, b, i}n_{e}(t)n_{i}
\end{equation}

Where $k_{r, b, i} [cm^3/s]$  and $n_i$ are the bimolecular recombination rate constant and species density $[1/cm^3]$ for the \emph{i}\textsuperscript{th} species, respectively. 

We are going to consider the situation where $G_{NE}$ is small compared to $G_{th}$, and make the following assumptions:

\begin{itemize}
\item $n_{K_+}(t) = n_{e}(t)$, where $n_{K_+}$ is the concentration of potassium ions. 
\item all other electron-capturing species $n_i$ are in thermal equilibrium.
\end{itemize}

This allows us to write

\begin{equation}
R = k_{r, b, K_{+}}n_{e}(t)^2 + \sum_{i \neq K+}^{}k_{r, m, i}n_{e}(t)
\end{equation}

The assumption that $n_i$ is in thermal equilibrium allows us to define the monomolecular recombination rate constant, $k_{r, m, i}$ 

\begin{equation}
\label{eq:monomolecular_recombination}
k_{r, m, i} = k_{r, b, i}n_{i}
\end{equation}


We write the $n_{e}(t)$ as being comprised of the equilibrium electron density, \(n_{e,0}\), and a perturbation, \(\Delta n_{e}(t)\).

\begin{equation}
  \label{eq:ne_pert}
n_{e}(t) = n_{e,0} + \Delta n_{e}(t)
\end{equation}

This allows us to write the recombination rate as

\begin{equation}
R = k_{r, b, K_{+}}(n_{e,0} + \Delta n_{e}(t))^2 + \sum_{i \neq K+}^{}k_{r, m, i}(n_{e,0} + \Delta n_{e}(t))
\end{equation}

expanding the square

\begin{equation}
R = k_{r, b, K_{+}}(n_{e,0}^2 + 2n_{e,0}\Delta n_{e}(t) + \Delta n_{e}(t)^2) + \sum_{i \neq K+}^{}k_{r, m, i}(n_{e,0} + \Delta n_{e}(t))
\end{equation}

We can group the time independent and time dependent terms together

\begin{align}
R = & k_{r, b, K_{+}}n_{e,0}^2 + \sum_{i \neq K+}^{}k_{r, m, i}n_{e,0} \nonumber \\
    & + 2k_{r, b, K_{+}}n_{e,0}\Delta n_{e}(t) + \sum_{i \neq K+}^{}k_{r, m, i} \Delta n_{e}(t) + k_{r, b, K_{+}}\Delta n_{e}(t)^2 
\end{align}

If we consider Equation \ref{eq:start_deqn} in the equilibrium case where $\frac{dn_{e}}{dt} = 0$ and $G_{NE} = 0$, we can see that the time independent terms are equal to $G_{th}$. Therefore the time independent terms will cancel out in Equation \ref{eq:start_deqn} and we can write

\begin{equation}
\frac{d\Delta n_{e}}{dt} = G_{NE}  -  2k_{r, b, K_{+}}n_{e,0}\Delta n_{e}(t) - \sum_{i \neq K+}^{}k_{r, m, i} \Delta n_{e}(t)- k_{r, b, K_{+}}\Delta n_{e}(t)^2
\end{equation}

Where we have used the fact that $n_{e,0}$ is constant in time. We also use the equation $n_{e,0} = n_{K_{+}}$ from our earlier assumption. 

\begin{equation}
\frac{d\Delta n_{e}}{dt} = G_{NE}   -   2k_{r, m, K_{+}}\Delta n_{e}(t) - \sum_{i \neq K+}^{}k_{r, m, i}\Delta n_{e}(t) - k_{r, b, K_{+}}\Delta n_{e}(t)^2
\end{equation}

The monomolecular recombination rates can be combined into an effective monomolecular recombination rate, $k_{r, m, eff}$. 

\begin{equation}
  k_{r, m, eff} = 2k_{r, m, K_{+}} + \sum_{i \neq K+}^{}k_{r, m, i}
\end{equation}

We then arrive at the final form of the differential equation for the time dependence of the electron density

\begin{equation}
  \label{eq:final_deqn}
\frac{d\Delta n_{e}}{dt} = G_{NE}   -   k_{r, m, eff}\Delta n_{e}(t) - k_{r, b, K_{+}}\Delta n_{e}(t)^2
\end{equation}



\subsection{MWS Kinetic fitting}

TODO: this probably can just be in the main text. 

For fitting of MWS data, we are interested in the time dependence after the laser pulse has ended. In this situation, we can assume that $G_{NE}$ is zero. Then Equation \ref{eq:final_deqn} becomes: 

\begin{equation}
\frac{d\Delta n_{e}}{dt} =   -   k_{r, m, eff}\Delta n_{e}(t)  - k_{r, b, K_{+}}\Delta n_{e}(t)^2
\end{equation}




\subsection{Photoionization Viability Assessment}

\subsubsection{Figure of Merit}

In this subsection we develop a model to assess the fundamental viability of photoionization as a method of increasing the power output of an MHD power system. We seek regimes where the power output from the MHD system will be increased by more than the power input of the nonequilibrium generation process. We define a figure of merit, $\gamma$ as the ratio of the increase in MHD power to the increase in power input.


\begin{equation}
  \label{eq:beta_def}
\gamma = \frac{dP_{MHD}}{dP_{in}} \Big|_{P_{in}=0} = \frac{dP_{MHD}}{dn_{e}}\frac{dn_{e}}{dP_{in}}\Big|_{P_{in}=0}
\end{equation}

Where we have indicated that we will evaluate $\gamma$ at $P_{in} = 0$. We have chosen to consider the effect of an infinitesimally small power input, as this will the best case scenario where the recombination rate is the lowest, as will be clarified below. A $\gamma > 1$ is possible as photoionization has the potential to significantly increase the efficiency of the MHD generator. 

For the MHD Power density we use the equation

\begin{equation}
  \label{eq:mhd_power}
P_{MHD} = \frac{1}{2} e \mu_e n_e u^2 B^2
\end{equation}

Where $e$ is the elementary charge, $\mu_e$ is the electron mobility, $n_e$ is the electron density, $u$ is the velocity of the gas, and $B$ is the magnetic field strength.  The determination of $\mu_e$ is discussed later. The derivative $\frac{dP_{MHD}}{dn_{e}}$ can be easily determined. 

For the equation $\frac{dn_{e}}{dP_{in}}$, we can use the equation \ref{eq:final_deqn} derived previously. We now develop an expression for $n_{e}(P_{in})$ and then take the derivative with respect to $P_{in}$.

We consider a gas subjected to a steady $G_{NE}$ and assuming the system is in quasi-equilibrium, $\frac{dn_{e}}{dt} = 0$. Note that $\ref{eq:ne_pert}$ can still be written with $\Delta n_e(t) = \Delta n_{e}$. (TODO need to think about how to rewrite general section with this.) With this, \ref{eq:final_deqn} which becomes

\begin{equation}
0 = G_{NE} - k_{r, m, eff}\Delta n_{e} - k_{r, b, K_{+}}\Delta n_{e}^2
\end{equation}

We now neglect the second term due to the fact that we are considering an infinitesimal power input, and therefore an infinitesimal $\Delta n_{e}$. For higher laser powers, the second term will increase the recombination rate. 

TODO: we could try and calculate the maximum power to which this would be valid. From our results we are entering the high-power regime where the second term is important with our pulsed laser. In any case, the argument of infinitesimal power input being best case scenario is valid and including the second term in general or exploring higher laser powers should be a lower proirity.

This leads to a simple expression for $\Delta n_{e}$

\begin{equation}
\Delta n_{e} = \frac{G_{NE}}{k_{r, m, eff}}
\end{equation}


We consider an arbitrary ionization mechanism that requires a power input per unit volume, \(P_{in}\), to generate a nonequilibrium generation rate \(G_{NE}\) . We can then define an ionization efficiency, \(\eta\) in the following way:

\begin{equation}
{G}_{NE} = \frac{\eta P_{In}}{E_{IP}}
\end{equation}


This equation is agnostic to the specific ionization mechanism, which would be encapsulated in \(\eta\). In this interpretation, \(G_{NE}\) represents the rate that electrons are lifted from the potential well of depth \(E_{IP}\). When \(\eta = 1\), the ionization process is perfectly efficient and for \(\eta < 1\) a fraction \(\eta - 1\) of the input energy is wasted. We note that \(\eta > 1\) is possible, particularly in the case of laser enhanced ionization where a laser excites an atom to a higher bound state, and a subsequent collisional excitation ionizes the atom. \(\eta\ \)on the order of 10 has been observed for metals in flames.{[}15{]} 

We can then write the electron density as a function of the power input

\begin{equation}
n_{e}(P_{in}) = \frac{\eta P_{In}}{E_{IP}k_{r, m, eff}}
\end{equation}

We can then take the derivative of $n_{e}(P_{in})$ with respect to $P_{in}$ to find $\frac{dn_{e}}{dP_{in}}$

\begin{equation}
\frac{dn_{e}}{dP_{in}} = \frac{\eta}{E_{IP}k_{r, m, eff}}
\end{equation}

Finally, we can combine these results to find $\gamma$

\begin{equation}
\gamma = \frac{1}{2} e \mu_e u^2 B^2 \frac{\eta}{E_{IP}k_{r, m, eff}}
\end{equation}

In the results presented in the main text we assume a velocity of 1000 m/s and a magnetic field strength of 5 Tesla. $E_{IP}$ is the ionization potential of potassium, 4.34 eV. In this work, we consider $\eta =1$, an $eta$ calculated from the energy absorbed by KOH (described later). For $k_{r, m, eff}$ we utilize equation \ref{eq:monomolecular_recombination}. The equilibrium species concentrations, $n_i$, and $\mu_e$ we use values calculated from from the chemical equilibrium model discussed in the next section. 




\subsubsection{Species concentration and mobility calculations}

To calculate $\mu_e$ and $n_i$ we model the combustion of seeded oxy-kerosene mixture with the Cantera python package. The input mass flow rates of $K_2CO_3$, $C_{12}H_{26}$, and $O_2$ are determined to achieve a specified equivalence ratio and potassium mass fraction (1\%). This mixture is initialized at a inlet temperature (300 K) then allowed to equilibrate at a fixed enthalpy and combustion chamber pressure (1e6 Pa).

We then model different potential gas areas of the MHD channel by equilibrating this gas mixture at various fixed temperature and pressures. Our Cantera model does not consider seed particle dissociation or formation. Figure 1 shows the calculated concentrations of the 10 highest-concentration species for our modeled oxy-fuel combustion gas mixture as function of temperature and atmospheric pressure.

We calculate $\mu_e$ of the gas mixture following the method outlined previously by Bedick et al.{[}14{]} Briefly, the collisional cross subsections for alkali vapors in flames. Then the electron mobility, \(\mu_{e}\), is calculated using the species-dependent collisional cross subsections of constituent gasses.

\includegraphics[width=4.09929in,height=3.95663in]{\repodir/final/figures/SI/output/Cantera_T.png}

Figure 1: a) Mole fractions of species with mole fraction of greater than 1e-3. The calculations are performed for oxy-kerosene combustion with \(\phi = 0.8\), and \(K_{wt} = 1\%\). b) Ionization fraction = {[}K+{]}/({[}K{]} + {[}K+{]}) and c) electrical conductivity.



\hypertarget{photoionization}{%
\subsubsection{UV Absorption}\label{photoionization}}



We apply our model to the specific case of UV photoionization. Photoionization, where potassium atoms are ionized by UV photons with energy above the potassium ionization potential, is a potential method of increasing the electrical conductivity of seeded combustion products. Photoionization is particularly attractive for contact targeting, due to the ability to direct and spatially confine light. Photoionization was considered as a method of increasing \(\sigma\) in MHD power systems early in MHD power research.{[}4{]} In that work, the potential of photoionization as an effective method of achieving sufficiently high electrical conductivity for MHD power was estimated. Those estimates were lacking in multiple ways:

\begin{itemize}
\item
  They considered blackbody radiation as the source of UV photons as UV laser technology was not developed yet. This is a highly inefficient way to generate UV photons. In addition, the use of lasers facilitates directional control of light, and advances in UV resistant fiber optics could potentially allow for precise application to specific areas.
\item
  The ionization of only a potassium vapor was considered, neglecting the absorption of light by other constituent gasses.
\item
  The calculations employed estimates of the potassium electron recombination rate, and experimental measurements of these reaction rates have been performed since then.
\end{itemize}

To apply our viability model to UV photoionization, we need to calculate the ionization efficiency for photoionization, \(\eta_{PI}\). We consider a fraction of the incident photons that are absorbed by potassium, hydroxide \(FA_{KOH}\). We assume a best case yield of KOH to electrons of 1. We define the efficiency as the ratio of the absorbed energy per incident photon to the ionization potential of potassium.


\begin{equation}
\eta_{PI} = \frac{E_{IP}}{E_{ph}}FA_{KOH} = \frac{E_{IP}}{E_{ph}}\frac{\delta}{\delta_{KOH}\ }FA\ \ 
\end{equation}

Where \(FA_{KOH}\) has been expressed in terms of the optical attenuation length of each individual species \(\delta_{i}\) and the total optical attenuation length of the gas mixture, \(\delta\), and the fraction of incident light absorbed by the entire gas mixture, FA. To understand the fundamental limit of photoionization we assume FA = 1. The case of FA for a practical system of an optical cavity with finite reflectivity mirrors is explored in the appendix.

\(\delta\) is calculated with the following equation:

\begin{equation}
\delta = \sum_{i}^{}\delta_{i} = \sum_{i}^{}{n_{i}Q_{i}}
\end{equation}

Where \(Q_{i}\) is the optical absorption cross subsection of the \emph{i}\textsuperscript{th} species. In this way \(\delta_{KOH} = n_{KOH}Q_{KOH}\)where the \(KOH\) stands for potassium.

To calculate \(\frac{\delta}{\delta_{KOH}\ }\) we utilize the compiled UV absorption cross subsection data outlined in \ref{fig:SI_UV_abscs} on the UV optical absorption cross subsections for the major species in the calculated gas mixtures.

Figure 6 shows\(\ \delta\), \(\delta_{KOH}\), and \(\frac{\delta}{\delta_{KOH}\ }\)for various UV wavelengths and temperatures.

\includegraphics[]{\repodir/modeling/viability/figure_panels/output/atten_length.png}

In the following work we analyze these gasses excited by 248 nm light. This is primarily as that is the longest wavelength of an excimer laser (KrF) that can still ionize potassium. However, the fraction of light absorbed by potassium is not particularly dependent on the wavelength, as can be seen in Figure 6. Although excimer lasers are pulsed we choose to analyze steady state illumination here. This is because for a given time-averaged laser intensity the time-averaged electron density will be higher with a steady state light source. The higher peak intensities of pulsed laser sources will create higher peak electron densities and a correspond higher rate of second order recombination. Therefore, considering steady state illumination corresponds to an upper bound on the electron density achievable with a given photoionization energy input. This assumption also allows for the calculation of equilibrium electron density by allowing for the equivalence of recombination rate and generation rate.

\subsubsection{UV Absorption Cross Sections}

We compiled data on UV absorption cross sections to determine potential ionization pathways. Data was collected within the relevant wavelength range for potassium photoionization (190-285 nm). The consolidated data is show in \ref{fig:SI_UV_abscs}. TODO: replace citations. 

For potassium we use compiled photoionization cross sections (Verner) for solar and blackbody radiation.{[}17{]} Regarding the other constituent gasses, we first rule out absorption by the atomic species H and O through the NIST database.{[}18{]} The UV absorption of the hydrogen molecule requires light with wavelength shorter than about 111 nm.{[}19{]} The OH radical similarly does not absorb UV light until wavelengths shorter than about 170 nm.{[}20{]} empirical formulas for the wavelength and temperature dependent UV cross sections exist for CO2{[}21{]}, H2O,{[}22{]} and O2.{[}22{]} KOH is dissociated by UV light in the 200-300 nm range and consequently has a relatively large absorption cross section. We extract the KOH absorption cross section data from Weng et all and extrapolate it in the range 190-200 nm.{[}23{]} Weng et al. observe a weak difference in absorption cross section for different flame conditions at 1400 and 1800 K so we do not include a temperature dependence to the KOH cross section. CO does absorb very weakly Cameron bands in the range of 170 -- 210 nm. We include CO absorption through the room temperature absorption measurements performed by Thompson et al.{[}24{]}


\begin{figure}[h]
    \includegraphics[width=0.8\linewidth]{\repodir/modeling/viability/figure_panels/output/abscs.png} 
    \caption{Optical absorption cross sections of species with significant cross section and UV absorption in the relevant range. Left: Fixed temperature wavelength dependence. Right: Fixed wavelength temperature dependence}
    \label{fig:SI_UV_abscs}
\end{figure}


% \clearpage

% \subsubsection{Supplemental PI Viability Results}

% \begin{figure}
% \centering
% \includegraphics[width=\textwidth]{\repodir/modeling/viability/figure_panels/output/beta_boundary_layer.png}
% \caption{Supplemental results}
% \end{figure}

% \begin{figure}
% \centering
% \includegraphics[width=\textwidth]{\repodir/final/figures/SI/output/Viability_rxn_compare.png}
% \caption{Supplemental results}
% \end{figure}


